\chapter{Introduction}\label{ch:intro}

\section{What \texorpdfstring{\nwsaom}{nwsaom} is}

\nwsaom{} estimates stochastic actor-oriented models (SAOMs) in Stata. A SAOM is a statistical model for network panel data: the same set of actors is observed at two or more points in time, and the model describes how the ties between them change from one observation to the next as the outcome of the decisions of the individual actors. The model was introduced by Snijders and co-workers; the standard introduction is Snijders, van de Bunt and Steglich~\cite{snijders2010}. Chapter~\ref{ch:model} summarises the model for readers who have not met it before.

\nwsaom{} is a Stata package of its own that builds on the \texttt{nwcommands} package. It works on networks that were declared with \texttt{nwset} (or created by one of the other \texttt{nw} commands), it reads actor attributes from the ordinary Stata dataset, and it returns its results in \texttt{e()} like any other Stata estimation command. The heavy computation, the simulation of the network evolution, runs in a compiled C plugin that is called from Mata. No R installation or other external statistical software is used when you fit a model.

\section{Relation to RSiena}

The reference implementation of SAOMs is the R package RSiena~\cite{ripley}, whose estimation function is \texttt{siena07} (called \texttt{siena} in the current RSiena manual). The other functions of RSiena that this manual refers to also have two names: \texttt{sienaGOF} is \texttt{test\_gof}, \texttt{sienaTimeTest} is \texttt{test\_time}, \texttt{sienaRI} is \texttt{interpret\_size}, \texttt{siena08} is \texttt{meta\_siena}, and \texttt{sienaBayes} of the package multiSiena is \texttt{multi\_siena}. The manual uses the older names, which are the ones in the literature. \nwsaom{} follows RSiena closely in its model, in its estimation algorithm, in the names of its effects and in the conventions that matter for the numbers (centring of covariates, the treatment of missing data, the scale of the rate parameters), so that RSiena's documentation and the published literature on SAOMs apply to \nwsaom{} as well.

\nwsaom{} is nevertheless an \emph{independent reimplementation}. It was written in Stata/Mata and C, using RSiena's published manual and, where the manual was not specific enough, RSiena's open-source code as the methodological reference. During development, RSiena (version 1.6.14, the current release) was used to certify the results: for the effects that were validated, the observed statistics (the targets of the method of moments) are identical to those of \texttt{siena07} and the estimates agree within simulation noise. The validation is documented effect by effect in the validation chapter. \nwsaom{} is not affiliated with, endorsed by or supported by the RSiena project.

Where the two programs differ in a way that matters to a user, the manual says so. The most important differences are these.
\begin{itemize}[nosep]
  \item Effects are chosen by Stata options (for example \opt{transtrip}, \opt{egox(alcohol1)}) and not by an effects table that is edited object by object.
  \item The goodness-of-fit test (\texttt{estat gof}) simulates new replicates at the fitted parameters instead of re-using networks stored during the fit; the test itself is the same Mahalanobis-distance test as in RSiena's \texttt{sienaGOF}.
  \item The option \opt{gwesp()} takes the decay parameter on the scale used in the \texttt{ergm} literature, which is $1/100$ of RSiena's parameter: RSiena's \texttt{gwespFF} with parameter 69 is \opt{gwesp(.69)}.
\end{itemize}

\section{What is covered}

\begin{itemize}[nosep]
  \item \textbf{Directed one-mode networks}, with two or more waves. With three or more waves the effects are pooled over the periods between consecutive waves and each period has its own rate.
  \item \textbf{Non-directed (symmetric) networks}, with RSiena's model types for non-directed relations (\opt{symtype()}).
  \item \textbf{Two-mode networks}, in which one set of nodes (the actors) chooses its ties to the other set (the events).
  \item \textbf{Behavior co-evolution}: a bounded integer behavior (an attitude, a count of cigarettes) evolves together with the network, so that selection and influence can be separated in one model; several behaviors are possible.
  \item \textbf{Several dependent networks} that evolve together on the same actors, with cross-network effects, through \texttt{nwsaom multinet}.
  \item \textbf{Several groups}: one model estimated jointly for several independent groups of actors (school classes, organizations) with the option \opt{groups()}, RSiena's \texttt{sienaGroupCreate} and \texttt{siena07} (Chapter~\ref{ch:groups}); or the separate fits combined afterwards with \texttt{nwsaom meta}.
  \item \textbf{Composition change}: actors that join or leave the network, either between waves (\opt{present()}) or at fractional times inside a period (\opt{joinat()}, \opt{leaveat()}).
  \item \textbf{Missing data}: missing ties (\opt{missnet()}), missing behavior values (\opt{missbeh()}) and missing covariate values; and \textbf{structural zeros and ones} (\opt{structural()}).
  \item \textbf{Rate covariates}: actor attributes that raise or lower an actor's rate of change (\opt{ratecov()}), and structural rate effects that depend on an actor's degree.
  \item \textbf{Covariates}: constant and time-varying actor covariates, and dyadic covariates.
  \item \textbf{Interactions} between effects, and the \textbf{endowment and creation} split of effects.
  \item \textbf{Other estimators}: maximum likelihood (\opt{ml}) and Bayesian estimation (\texttt{nwsaom bayes}) for a directed one-mode network, with or without one co-evolving behavior (Chapter~\ref{ch:estimation}).
  \item \textbf{Postestimation and tools}: goodness of fit (\texttt{estat gof}), micro effects on macro structure (\texttt{estat mems}), score-type tests (\texttt{estat score}), the time-heterogeneity test (\texttt{estat timetest}), the relative importance of effects (\texttt{estat importance}), simulation at chosen parameters (\texttt{nwsaom simulate}), a description of the data before a fit (\texttt{nwsaom describe}), and a meta-analysis of the same model fitted to several networks (\texttt{nwsaom meta}), the counterpart of RSiena's \texttt{siena08}.
\end{itemize}
The effects library follows RSiena's table of effects. Of the \nTotal{} rows of the effects table of RSiena 1.6.6 (\texttt{getEffects()} over twelve data configurations), about \pctImpl\% of the rows in scope are implemented, \nImpl{} in all; (The counts refer to RSiena 1.6.6; the effects that RSiena 1.6.7 to 1.6.14 added or changed, for example \texttt{WXZX} and the two-mode \texttt{toU}, are implemented and compared with RSiena 1.6.14 but are not part of these counts, and the counts change with RSiena releases.)  Chapter~\ref{ch:effects} and Appendix~\ref{app:effects} give the details.

\paragraph{What RSiena has that \nwsaom{} does not.} \nwsaom{} does not cover all of RSiena. Not implemented are the generalised method of moments estimator (31 statistics), the settings model (12 effects), a hierarchical Bayesian analysis of several groups, and the continuous (stochastic differential equation) behavior (80 rows, not planned). Inside \texttt{nwsaom multinet} there is no composition change, no missing ties and no \texttt{estat gof}, \texttt{estat importance} or \texttt{estat simulate}. Maximum likelihood and Bayesian estimation do not cover non-directed or two-mode networks, several networks, missing data or composition change. The joint estimation of several groups (\opt{groups()}) does not yet cover maximum likelihood and Bayesian estimation, \texttt{estat gof}, \texttt{estat simulate} and \texttt{estat importance}, two-mode and non-directed networks, several networks or behaviors, and a few effects (Section~\ref{sec:groups-limits}). Chapter~\ref{ch:missing} lists every missing effect with its reason.

Binary ties only are supported. As RSiena, \nwsaom{} judges a network by its tie values: every tie must have the value 1, 10 or 11. A network with other tie values (real shares, counts) stops the command with an error and is never silently dichotomised; a network that is flagged as valued by \texttt{nwcommands} but whose ties are all 1 is accepted. The values 10 and 11 are RSiena's codes for structural zeros and ones; \nwsaom{} detects them in the wave networks (Section~\ref{sec:structural}). All waves of a model must have the same number of nodes; the actors are identified by their position, so node $i$ must be the same actor in every wave.

\section{Installation and requirements}

\nwsaom{} needs the \texttt{nwcommands} core package, which provides the network objects, \texttt{nwset} and the data-handling commands. Install the core first and then \nwsaom{}:
\begin{lstlisting}
. net install nwcommands, from(https://nwcommands.org/net/) replace
. nwinstall, all
. nwinstall, modelpkg(nwsaom)
\end{lstlisting}
The first line installs the core, \texttt{nwinstall, all} installs its help files, dialog boxes and menu, and \texttt{nwinstall, modelpkg(nwsaom)} installs \nwsaom{} from its repository. To install the core and all four model packages (\texttt{nwsaom}, \texttt{nwergm}, \texttt{nwdynam}, \texttt{nwrem}) in one step, use
\begin{lstlisting}
. net install nwcommands, from(https://nwcommands.org/net/) replace
. nwinstall, everything
\end{lstlisting}
\nwsaom{} can also be installed by itself, once the core is there:
\begin{lstlisting}
. net install nwsaom, from(https://raw.githubusercontent.com/thomas-u-grund/nwsaom/main/) replace
\end{lstlisting}
The core must be version 3.6.0 or later (the listings of this manual were checked with core 3.9.2). \nwsaom{} checks this when it starts, and also whether the Mata libraries of the core and of \nwsaom{} were built together, and says how to reinstall if they were not. (The Mata library of \nwsaom{} contains the class definitions of the core that it was compiled against; a core of a different build would otherwise stop Mata with the error \texttt{class compiled at different times}.) The Mata libraries are indexed automatically, so that a fresh installation works in a new Stata session without further steps. Installing or updating \texttt{nwcommands} later never replaces \nwsaom{}, because the core contains no file with the name of a \nwsaom{} command. Alternatively, add a clone of the \nwsaom{} repository to the ado-path with \opt{adopath + "/path/to/nwsaom"}; the tests of the repository find the core through the environment variable \texttt{NWCORE}.

\paragraph{How to cite.} The manual is the citable reference for the package; it is archived at Zenodo and has a DOI. Please cite it as: Grund, T.~U. (2026). \emph{nwsaom: Stochastic actor-oriented models in Stata. Manual for version \manualseries.} Zenodo. \url{https://doi.org/\manualdoi}. The label is the series of the manual (it holds for all versions \manualseries.x); the DOI is the one of all versions of the record. Please also report the exact version that produced your results, as printed by \texttt{nwsaom version} (this manual: \pkgversion; Section~\ref{sec:postest}).

\paragraph{Versions.} \texttt{nwsaom version} reports the installed release of \nwsaom{}, the version of the \texttt{nwcommands} core, the native plugin in use, the version of RSiena that the results were compared with and the license. \texttt{nwversion} lists the core and all installed model packages. \texttt{net install nwsaom} installs the state of the main branch. A specific release is installed from its tag, for example
\begin{lstlisting}
. net install nwsaom, from(https://raw.githubusercontent.com/thomas-u-grund/nwsaom/v1.2.0/) replace
\end{lstlisting}

\paragraph{What changed in the 1.2 series.} This manual describes the behavior of version 1.2; the points below are the changes since the previous archived manual (labelled 1.1), in case you re-run old analyses. Results of models that were possible before are unchanged.
\begin{itemize}[nosep]
  \item \emph{Joint estimation of several groups:} the options \opt{groups()} and \opt{gwaves()}, with the rates numbered as RSiena numbers the periods of the combined series, the stored results \texttt{e(rates\_group)}, \texttt{e(rates\_period)}, \texttt{e(groupn)}, \texttt{e(groupwaves)}, \texttt{e(ngroups)} and \texttt{e(groupvar)}, and \texttt{estat timetest} and \texttt{estat score} after such a fit (Chapter~\ref{ch:groups}).
  \item \emph{Two-mode networks:} \texttt{nwsaom describe} works for a two-mode network (the actors, events and attendances; before, it counted the cells within the modes and reported an undirected network of all nodes) and has the option \opt{actormode()}; \texttt{estat simulate} and \texttt{nwsaom simulate} work after and for a two-mode model, and \opt{generate()} saves two-mode networks (Section~\ref{sec:describe} and Section~\ref{sec:sim}). A \opt{theta()} matrix of the wrong shape gives a message.
  \item \emph{\texttt{nwsaom meta, rates}} used standard errors that did not exist, so every rate had a missing standard error and was silently excluded; it now uses \texttt{e(rates\_se)}. The result matrices \texttt{r(theta)} and \texttt{r(se)} now carry the names of the stored estimates (rows) and of the effects (columns).
  \item \emph{A message:} the refusal of \opt{reciprocity} for a non-directed network no longer prints a second line.
  \item \emph{Speed} (version 1.1.1): estimation with the compiled plugin is faster, most for two-mode networks; the results are identical.
\end{itemize}

\paragraph{What changed in the 1.1 series.} The points below are the changes between the manual labelled 1.0.3 and the manual labelled 1.1.
\begin{itemize}[nosep]
  \item \emph{Tie values are judged by their values.} A network with any tie value other than 1 (real shares, counts, 10 or 11) stops \nwsaom{}, \texttt{nwsaom multiplex} and \texttt{nwsaom multinet} with an error; before, a network of shares that had been saved as not valued was turned into ties silently. A network flagged as valued whose ties are all 1 is accepted.
  \item \emph{A visible warning for the Mata fallback} (Section~\ref{sec:fallbackwarn}).
  \item \emph{Installation and requirements:} \texttt{nwinstall} installs the package; the core must be version 3.6.0 or later; Stata 17 or later; the Mata libraries are built with Stata 17 (libraries of earlier releases were built with Stata 18 and could not be used in Stata 17); a fresh installation works without \texttt{mata mlib index}; a mismatch between the Mata libraries of the core and of \nwsaom{} gives a clear message instead of the error r(3021).
  \item \emph{Convergence safeguards} for rate effects and behavior rates (Chapter~\ref{ch:estimation}).
  \item \emph{Messages and dialogs:} the messages about a plugin that is too old now say that a newer plugin is needed; \texttt{nwsaom version} prints ``plugin build'' for the plugin; every option and effect can be set from a dialog box (\texttt{db nwsaom}); the version is shown at the bottom of each help file.
\end{itemize}

\paragraph{License.} \nwsaom{} is free software under the GNU General Public License, version~3. It is derived from RSiena (Ripley, Snijders, Boda, V\"or\"os and others; GPL-2 or GPL-3) and, for the Bayesian estimator, from the algorithm of the multiSiena package (Snijders, Koskinen and Ripley; GPL-3). The source code is public; the file \texttt{NOTICE.md} of the repository lists the sources and the version of RSiena that each part was compared with.

\paragraph{Software.} \nwsaom{} needs Stata~17 or later and was tested with Stata~17 and Stata~18 (Stata/BE is enough); older versions of Stata have not been tested. The Mata libraries are built with Stata~17, the oldest tested version, because a library built by a newer Stata is ignored by older ones. The listings of this manual were produced with Stata/BE~18. The edition of Stata determines how large a network can be when it is passed as a matrix (Section~\ref{sec:limits}).

\paragraph{The native plugin.} The simulator is the compiled plugin \texttt{saom\_sim}. The repository contains builds for macOS (\opt{saom\_sim\_macos.plugin}), for Linux (\opt{saom\_sim\_unix.plugin}) and for Windows (\opt{saom\_sim\_windows.plugin}). The macOS build is compiled on the development machine; the Linux and Windows builds are produced by continuous integration and can be older than the \texttt{.ado} files. A plugin that is older than a model needs is detected (the plugin reports its version), and the model is then either fitted by the slower Mata simulator or, for the effects that exist only in the plugin, refused with an error message. Whenever a fit runs on the Mata simulator, \nwsaom{} prints one warning line before the estimation starts and says why (Section~\ref{sec:fallbackwarn}); you can also look at \texttt{e(engine)} afterwards.

\section{How to read this manual}

If you know SAOMs from RSiena, read Chapter~\ref{ch:start} for the syntax of a first model and use Chapter~\ref{ch:options} as a dictionary from RSiena concepts to \nwsaom{} options. If you are new to SAOMs, read Chapter~\ref{ch:model} first; it is short and uses only the equations that you need to interpret the output.

\begin{description}[style=nextline,nosep]
  \item[Chapter~\ref{ch:model}] The model in brief: rate function, objective function, ministeps, the method of moments and its three phases, conditional and unconditional estimation, convergence.
  \item[Chapter~\ref{ch:start}] How the data must look, a complete first analysis with real output, and the results that \nwsaom{} stores.
  \item[Chapter~\ref{ch:options}] Reference of all options of \texttt{nwsaom}, organised by topic.
  \item[Chapter~\ref{ch:effects}] The effects: what each statistic counts, its RSiena name and the option that selects it.
  \item[Later chapters] Behavior co-evolution, the network types, rates and covariate data, estimation practice, validation against RSiena, and composition change and missing data.
  \item[Chapter~\ref{ch:groups}] Several groups: joint estimation with \opt{groups()}, when to use it instead of the meta-analysis, and what is not yet available with groups.
  \item[Chapter~\ref{ch:technical}] Technical notes: the plugin, threads, size limits, performance and the file layout of the source.
\end{description}

\paragraph{Typographic conventions.} Stata commands and options are set in \texttt{typewriter}. Output shown in listings was produced by the commands above it with the version of \nwsaom{} described in this manual; the random-number seed is always given, so that the output can be reproduced. Simulation-based estimates change in the last digits when the model or the seed changes, and, for a given seed, they do not depend on the number of processor cores.
